\documentclass[aspectratio=169,10pt]{beamer}
\usetheme{default}
\usecolortheme{seahorse}
\usepackage[utf8]{inputenc}
\usepackage[T1]{fontenc}
\usepackage{lmodern}
\usepackage{graphicx}
\usepackage{booktabs}
\usepackage{amsmath}
\usepackage{xcolor}

\graphicspath{{../figures/}}
\setbeamertemplate{navigation symbols}{}
\setbeamertemplate{footline}[frame number]
\setbeamerfont{frametitle}{size=\large}
\setbeamercolor{frametitle}{fg=blue!35!black}
\setbeamercolor{structure}{fg=blue!35!black}
\newcommand{\ns}{\,\mathrm{ns}}
\newcommand{\hl}[1]{\textcolor{orange!75!black}{\textbf{#1}}}

\title[Flash timing]{The $\gamma$-flash time base of the n\_TOF EAR2 X17 campaign}
\subtitle{Calibration from the seven runs with the SiPM-wall divert disabled}
\author{n\_TOF X17 / DREAM analysis}
\date{2026-07-28}

\begin{document}

\begin{frame}[plain]
  \titlepage
\end{frame}

% ------------------------------------------------------------------
\begin{frame}{The problem, in one slide}
\begin{columns}[T]
\begin{column}{0.54\textwidth}
\begin{itemize}\setlength\itemsep{0.5em}
  \item The SiPM walls are \hl{blanked} around the $\gamma$-flash
        ($\sim1\,\mu$s gate, N1081B M6 sec.~C).
  \item So the PSA flash finder tags either
        \begin{itemize}
          \item the \textbf{gate transient} ($-1190\ns$ before 07-22,
                $-375\ns$ after), or
          \item a \textbf{clamped leak} at the true flash time,
        \end{itemize}
        and it \hl{switches between the two from bunch to bunch}.
  \item $\Rightarrow$ no per-run offset can repair the stored
        \texttt{tflash}. It has to be replaced.
\end{itemize}
\vspace{0.6em}
\begin{block}{The fix}
\[
  t_\mathrm{flash} = t_\mathrm{PKUP} + C,\qquad C\simeq-1719\ns
\]
The pickup is never gated and its finder never fails.
\end{block}
\end{column}
\begin{column}{0.44\textwidth}
  \includegraphics[width=\linewidth]{08_plastic_anatomy.png}\\
  {\tiny The flash region. Right-hand panel: the wall shows the gate transient
   at $-2100\ns$ \emph{and} the leak at $-1750\ns$.}
\end{column}
\end{columns}
\end{frame}

% ------------------------------------------------------------------
\begin{frame}{Finding the runs: a scan of the whole campaign}
\begin{columns}[T]
\begin{column}{0.5\textwidth}
306 processed runs scanned (224264--224587):
\vspace{0.4em}
\begin{itemize}
  \item \hl{7 runs} with the divert OFF
  \item 154 runs with it ON
  \item 145 runs with no wall trees at all (before 224345) --- unverifiable
        without a tape reprocess
\end{itemize}
\vspace{0.6em}
Signature: all four walls tag the same time within $\pm6\ns$, and the flash peak
is \hl{$\sim$28\,500 ADC} instead of the $\sim$850 ADC gate transient.
\vspace{0.5em}

Two independent epochs, five days apart --- which is what makes the
cross-checks possible.
\end{column}
\begin{column}{0.5\textwidth}
\small
\begin{tabular}{lrl}
\toprule
run & bunches & intensity \\
\midrule
224356 & 3412 & dedicated \\
\textbf{224357} & 3286 & 47/53 ded./par. \\
224358 & 3269 & mixed \\
224359 & 3426 & dedicated \\
224360 & 549 & dedicated \\
\midrule
224464 & 2716 & mostly ded. \\
224466 & 2073 & mostly ded. \\
\bottomrule
\end{tabular}
\vspace{0.8em}

{\footnotesize 07-11 group: N1081B M6 offline during the switch outage.\\
07-16 pair: during the wall/plastic recalibration.}
\end{column}
\end{columns}
\end{frame}

% ------------------------------------------------------------------
\begin{frame}{Per-channel constants}
\begin{center}
\includegraphics[width=0.85\textwidth]{01_per_channel_arrival.png}
\end{center}
\vspace{-0.4em}
{\footnotesize
Within an epoch every channel reproduces to \hl{$0.6\ns$}. The
$07\text{-}11\rightarrow07\text{-}16$ shift is \hl{real and per-channel}
(WALB $-0.6\ns$, WALD $-6.9\ns$): the wall HV/threshold equalisation of 07-15/16.
After it, the four walls agree to $0.7\ns$ --- and that is the state that held
for the rest of the campaign.}
\end{frame}

% ------------------------------------------------------------------
\begin{frame}{The calibration}
\begin{columns}[T]
\begin{column}{0.52\textwidth}
\small
\begin{tabular}{lrrr}
\toprule
wall & all 7 & 07-11 & \textbf{07-16} \\
\midrule
WALA & $-1717.5$ & $-1716.7$ & $\mathbf{-1719.5}$ \\
WALB & $-1719.3$ & $-1719.1$ & $\mathbf{-1719.7}$ \\
WALC & $-1715.5$ & $-1714.0$ & $\mathbf{-1719.2}$ \\
WALD & $-1714.1$ & $-1712.2$ & $\mathbf{-1719.0}$ \\
\midrule
all 32 & $-1716.6$ & $-1715.5$ & $\mathbf{-1719.3}$ \\
\bottomrule
\end{tabular}
\vspace{0.8em}

{\footnotesize Recommended: the \textbf{07-16} column, per channel, for runs
$\ge224400$.}
\end{column}
\begin{column}{0.48\textwidth}
\small
\begin{tabular}{lr}
\toprule
uncertainty & [ns] \\
\midrule
run-to-run, within epoch & 0.5 \\
time base over campaign (rms) & 0.68 \\
transport to end of campaign & 1.2 \\
07-11 vs 07-16 epoch & 3.9 \\
intensity, par.\ $\to$ ded. & 5.0 \\
\midrule
single bunch, one channel & 3.5 \\
single bunch, 32-ch average & 2.5 \\
\bottomrule
\end{tabular}
\end{column}
\end{columns}
\end{frame}

% ------------------------------------------------------------------
\begin{frame}{Timing resolution: what is in the 3.5 ns}
\begin{center}
\includegraphics[width=0.78\textwidth]{04_distributions.png}
\end{center}
\vspace{-0.3em}
\begin{columns}[T]
\begin{column}{0.5\textwidth}
{\footnotesize
Decomposing the 32-channel matrix, per bunch:
\begin{itemize}
  \item single channel vs pickup: $3.5\ns$
  \item common to all 32 channels: $2.5\ns$
  \item independent per channel: $2.5\ns$
\end{itemize}}
\end{column}
\begin{column}{0.5\textwidth}
{\footnotesize
The independent part averages away over the array; the common part does not.
The pickup contributes only ${\sim}1\ns$, so the common $2.5\ns$ is a
\hl{genuine bunch-to-bunch variation of the flash arrival}.\\[0.3em]
07-16 runs are better than 07-11 ($2.7$ vs $3.9\ns$) --- again the equalisation.}
\end{column}
\end{columns}
\end{frame}

% ------------------------------------------------------------------
\begin{frame}{The reference: why the pickup}
\begin{columns}[T]
\begin{column}{0.5\textwidth}
\begin{itemize}\setlength\itemsep{0.45em}
  \item per-bunch spread \hl{$0.83\ns$} (dedicated), $1.42\ns$ (parasitic)
  \item usable in \hl{98.4\,\%} of bunches\\
        {\footnotesize (the 1.6\,\% failures are all parasitic)}
  \item amplitude strictly \hl{linear in protons}: ratio $1.97$ for a factor 2
  \item timing moves by only $-0.45\ns$ between intensity classes
  \item its flash finder \hl{never fails} --- unlike WAL/PSS/LIQ
\end{itemize}
\end{column}
\begin{column}{0.5\textwidth}
\includegraphics[width=\linewidth]{03_stability.png}\\
{\footnotesize Within a full three-hour run the flash time is flat to
$0.25$--$0.39\ns$ bin-to-bin, $1.4\ns$ peak-to-peak.}
\end{column}
\end{columns}
\end{frame}

% ------------------------------------------------------------------
\begin{frame}{Question 1: is the intensity dependence real?}
\begin{center}
\includegraphics[width=0.95\textwidth]{06_intensity_mechanism.png}
\end{center}
\vspace{-0.3em}
{\footnotesize
The wall constant moves by $-5.0\ns$ from parasitic to dedicated pulses. But:
\begin{itemize}
  \item the shift is \hl{detector-dependent} --- $0$ (pickup), $-5$ (walls),
        $-3\ldots-6$ (plastics), $-16\ns$ (silicon). A real change of the flash
        arrival would be \emph{common to all}.
  \item only the pickup is linear in protons; everything else is saturated
        (amplitude ratio $\simeq1$)
  \item the shift follows the \hl{risetime change}: walls $20.5\to15.4\ns$,
        silicon $98\to56\ns$
\end{itemize}}
\end{frame}

% ------------------------------------------------------------------
\begin{frame}{Question 1: and it is not classical amplitude walk}
\begin{columns}[T]
\begin{column}{0.5\textwidth}
\begin{center}
\includegraphics[width=\linewidth]{07_intensity_prediction.png}
\end{center}
\end{column}
\begin{column}{0.5\textwidth}
\begin{itemize}\setlength\itemsep{0.5em}
  \item Measure the amplitude$\to$time slope \emph{within} one intensity class.
  \item Extrapolate it to the observed amplitude change \emph{between} classes.
  \item Prediction: $+0.3\ns$ (walls, \hl{wrong sign}), $<0.4\ns$ elsewhere.
  \item Observed: $-5$ to $-16\ns$.
\end{itemize}
\vspace{0.5em}
\begin{block}{Mechanism}
{\small The front end saturates: a bigger flash cannot make a taller pulse, so it
reaches the clipping level \hl{earlier} and stays there longer
(FWHM $700\to2000\ns$). Any edge estimator fires early.}
\end{block}
\vspace{0.3em}
{\footnotesize $\Rightarrow$ carry it in the calibration; do \emph{not} read it
as physics.}
\end{column}
\end{columns}
\end{frame}

% ------------------------------------------------------------------
\begin{frame}{Question 2: is the time base stable across the campaign?}
\begin{center}
\includegraphics[width=0.63\textwidth]{11_timeseries.png}
\end{center}
\vspace{-0.4em}
{\footnotesize
LIQ and PSS are \hl{never gated}, so they see the flash in every run --- the only
way to test stability where the walls cannot be measured.}
\end{frame}

% ------------------------------------------------------------------
\begin{frame}{Question 2: the answer}
\begin{columns}[T]
\begin{column}{0.5\textwidth}
\textbf{Liquids: stable}
\begin{itemize}
  \item $\sigma_\mathrm{run} = 0.73,\,0.73,\,0.77,\,1.48\ns$ (A--D) over 36--39
        runs, 224400--224584
  \item the four cells \hl{move together}: within-run spread of deviations
        $0.22\ns$, LIQA--LIQC correlation $+0.92$
  \item common-mode wander \hl{$0.68\ns$ rms}, worst single run $-2.2\ns$
        (224531)
  \item the 07-16 divert-off runs sit $+0.9\ldots+1.2\ns$ from the late-campaign
        mean $\Rightarrow$ \hl{transport good to $1.2\ns$}
\end{itemize}
\end{column}
\begin{column}{0.5\textwidth}
\textbf{Independent check in gated runs}
\begin{itemize}
  \item the clamped wall leak lands within \hl{$6.2\ns$} of the calibration
        (worst case $9.2\ns$), on all 32 channels
  \item in both 224362 (07-11 era) and 224572 (07-26, 15 days later)
  \item the gate transient sits at $-1190$ and $-375\ns$ --- reproducing the
        $+800\ns$ delay change of 07-22
\end{itemize}
\vspace{0.5em}
\begin{block}{}
{\small The wall calibration measured on 07-16 carries to the end of the
campaign at the ${\sim}1\ns$ level.}
\end{block}
\end{column}
\end{columns}
\end{frame}

% ------------------------------------------------------------------
\begin{frame}{Question 3: the plastics --- anatomy}
\begin{center}
\includegraphics[width=0.95\textwidth]{08_plastic_anatomy.png}
\end{center}
\vspace{-0.2em}
{\footnotesize
The wall gives \hl{one clean pulse}. The plastic response is spread over
${\sim}100\ns$ and sits on a continuum of fragments: \hl{25--30 hits per bunch}
in the $\pm1\,\mu$s window, against 12 for the walls.}
\end{frame}

% ------------------------------------------------------------------
\begin{frame}{Question 3: the plastics --- three failure modes}
\small
\begin{tabular}{llrrrrrr}
\toprule
epoch & family & found & $\sigma$ [ns] & amplitude & sat.\ flag & FWHM & $n_\mathrm{hit}$/b \\
\midrule
2026-07-11 & WAL & 100\,\% & 4.1 & 28521 & 0\,\% & 1527 & 12.4 \\
2026-07-11 & PSS & \hl{67\,\%} & \hl{19.1} & 23193 & 0\,\% & 11 & 38.2 \\
2026-07-16 & WAL & 100\,\% & 2.3 & 28560 & 0\,\% & 2602 & 12.5 \\
2026-07-16 & PSS & 64\,\% & 5.8 & 34571 & 0\,\% & 43 & 26.0 \\
2026-07-16 & LIQ & \textbf{100\,\%} & 4.1 & \hl{63931} & \hl{80\,\%} & 103 & 43.6 \\
2026-07-26 & WAL & 100\,\% & 8.3 & 19650 & 6\,\% & 57 & 11.7 \\
2026-07-26 & PSS & \hl{42\,\%} & 2.5 & \hl{55114} & 27\,\% & 76 & 24.3 \\
2026-07-26 & LIQ & \textbf{100\,\%} & 4.5 & \hl{63921} & \hl{79\,\%} & 102 & 44.1 \\
\bottomrule
\end{tabular}
\vspace{0.5em}

\begin{enumerate}\setlength\itemsep{0.3em}
  \item \textbf{Yield.} A flash hit exists in only 64--67\,\% of bunches,
        \hl{42\,\%} after the FIFO fan-in. In the missing bunches the PSA
        reconstructs \emph{no hit at all} before ${\sim}-1600\ns$ --- verified
        not to be a selection artefact.
  \item \textbf{Amplitude-dependent bias} up to $17\ns$, $\sigma=7$--$15\ns$
        (pre-FIFO).
  \item \textbf{Instability}: the constant moves tens of ns across the campaign
        ($\sigma_\mathrm{run}=17$--$31\ns$).
\end{enumerate}
\end{frame}

% ------------------------------------------------------------------
\begin{frame}{Question 3: it is not ``saturation breaks the PSA''}
\begin{columns}[T]
\begin{column}{0.52\textwidth}
The \hl{liquid scintillators are the control}: same digitiser, same PSA,
different front end.
\vspace{0.5em}
\begin{itemize}\setlength\itemsep{0.4em}
  \item LIQ flash hits are \hl{more} saturated than the plastics':
        sat.\ flag \textbf{79--80\,\%} vs 27\,\%, amplitude 63\,900 vs 55\,100
  \item highest hit multiplicity of all: \textbf{44}/bunch
  \item and yet the PSA finds the flash in \textbf{100\,\%} of bunches, at
        $\sigma = 4.1$--$4.5\ns$
\end{itemize}
\end{column}
\begin{column}{0.48\textwidth}
\begin{block}{So the plastic drop-out is\ldots}
{\small \hl{specific to the plastic channel} --- its PSA configuration and/or its
pulse shape (PSS flash FWHM 43--76\,ns and enormous area) --- \textbf{not} a
generic consequence of a railed pulse.}
\end{block}
\vspace{0.6em}
{\small $\Rightarrow$ plausibly \hl{fixable in the \texttt{UserInput}}, which is
exactly what the reprocessing effort is for.}
\end{column}
\end{columns}
\end{frame}

% ------------------------------------------------------------------
\begin{frame}{Question 3: precision vs pulse size, and per channel}
\begin{columns}[T]
\begin{column}{0.5\textwidth}
\includegraphics[width=\linewidth]{09_plastic_vs_amp.png}
\end{column}
\begin{column}{0.5\textwidth}
\includegraphics[width=\linewidth]{10_plastic_channels.png}
\end{column}
\end{columns}
\vspace{0.3em}
{\footnotesize
The two regimes trade off: pre-FIFO the hit is usually found but badly timed;
post-FIFO it is precisely timed ($2.4$--$3.4\ns$) but usually missing.
PSSB\,ch2 is defective throughout (amplitude 9\,778 vs ${\sim}30\,000$).
$\Rightarrow$ take plastic constants \hl{per run}, never per epoch.}
\end{frame}

% ------------------------------------------------------------------
\begin{frame}{Update: the plastics are fixed by the reprocessing}
\begin{columns}[T]
\begin{column}{0.5\textwidth}
The parallel UserInput work reprocessed run 224572. Same measurement on
\texttt{v4\_walshapes}:
\vspace{0.5em}

\small
\begin{tabular}{lrr}
\toprule
 & official & reprocessed \\
\midrule
PSS flash hit found & 42\,\% & \hl{100\,\%} \\
per-bunch $\sigma$ & 2.5\,ns & 3.0--4.6\,ns \\
LIQ found & 100\,\% & 100\,\% \\
\bottomrule
\end{tabular}
\vspace{0.8em}

{\footnotesize The single change responsible is the \hl{G-FLASH threshold}
($50\rightarrow2000/10^4$): \texttt{v1\_flash} alone, elimination untouched,
already gives 100\,\%.}
\end{column}
\begin{column}{0.5\textwidth}
\begin{block}{So the missing hits were\ldots}
{\small a \hl{downstream symptom of the flash-finder bug}, not an independent
plastic defect. With the flash mis-located at ${\sim}130\ns$ (fitting noise),
the pulse recognition never reconstructs the real pulse at $11.63\,\mu$s.}
\end{block}
\vspace{0.5em}
{\footnotesize Elimination is \emph{not} the gate: widening it
(\texttt{v2\_elim}) changes the yield not at all, and the surviving flash hits
have area/amp $\simeq72$--$76$, above even the widened window.}
\vspace{0.5em}

{\footnotesize $\Rightarrow$ \S\,``plastics unusable'' applies to the
\emph{official} files only.}
\end{column}
\end{columns}
\end{frame}

% ------------------------------------------------------------------
\begin{frame}{Independent cross-check: the two routes agree}
\small
\begin{columns}[T]
\begin{column}{0.56\textwidth}
\begin{tabular}{lrrr}
\toprule
tree & reprocessed & this work & diff \\
\midrule
WALA & $-1717.82$ & $-1719.46$ & $+1.64$ \\
WALB & $-1718.07$ & $-1719.69$ & $+1.62$ \\
WALC & $-1718.05$ & $-1719.23$ & $+1.18$ \\
WALD & $-1722.29$ & $-1719.01$ & $-3.28$ \\
\multicolumn{3}{r}{rms} & \hl{2.07} \\
\midrule
LIQA & $-1708.38$ & $-1708.22$ & $-0.16$ \\
LIQB & $-1710.78$ & $-1710.30$ & $-0.48$ \\
LIQC & $-1695.73$ & $-1695.56$ & $-0.17$ \\
LIQD & $-1701.27$ & $-1701.56$ & $+0.29$ \\
\multicolumn{3}{r}{rms} & \hl{0.27} \\
\bottomrule
\end{tabular}
\end{column}
\begin{column}{0.44\textwidth}
{\footnotesize
A \hl{corrected flash finder} run on 224572 --- a \emph{gated} run, ten days
after the nearest calibration run, where the walls see only the clamped leak.
\\[0.5em]
Walls agree to \hl{2.1\,ns} rms, liquids to \hl{0.27\,ns} rms --- two different
processings of the same run agreeing to a quarter of a nanosecond.
\\[0.5em]
The two routes share no assumption beyond the pickup itself.}
\end{column}
\end{columns}
\end{frame}

% ------------------------------------------------------------------
\begin{frame}{What to do}
\begin{enumerate}\setlength\itemsep{0.55em}
  \item \textbf{Walls}: per-channel \texttt{C\_2026\_07\_16} for runs
        $\ge224400$, \texttt{C\_2026\_07\_11} before. Add $-5\ns$ if the pulse
        type differs from the calibration one.
  \item \textbf{Liquids}: per-run value or campaign mean --- they agree to
        ${\sim}1\ns$.
  \item \textbf{Plastics}: per-run value only.
  \item \hl{Never} use the stored \texttt{tflash} of WAL/PSS/LIQ \emph{from the
        official files}; in reprocessed files it is sound, but PKUP\,$+\,C$ is
        still the recommended time base.
  \item If a bunch has no usable pickup pulse (1.6\,\%, all parasitic), drop the
        bunch.
\end{enumerate}
\vspace{0.8em}
\begin{block}{Limitations}
{\small Amplitude is unusable in these runs (front end saturated) --- energy
scale still from the $^{88}$Y runs 224476--479. The $3.9\ns$ epoch shift shows
wall timing \emph{can} move; nothing rules out a further move after 07-16 beyond
the liquid evidence. Runs before 224345 unverified. Raw data of all seven runs
has aged off EOS disk.}
\end{block}
\end{frame}

% ------------------------------------------------------------------
\begin{frame}{Where everything lives}
\small
\begin{tabular}{ll}
\toprule
\texttt{data/flash\_timing\_calibration.json} & the calibration, per channel and epoch \\
\texttt{data/per\_channel\_flash\_timing.csv} & 280 measurements behind it \\
\texttt{data/divert\_state\_by\_run.csv} & the 306-run divert scan \\
\texttt{data/plastic\_liq\_flash\_by\_run.csv} & the time-base monitor \\
\texttt{data/plastic\_pathology.csv} & the plastic study \\
\texttt{data/resolution\_decomposition.json} & jitter split, per run \\
\texttt{figures/}, \texttt{scripts/}, \texttt{latex/} & plots, code, this deck \\
\bottomrule
\end{tabular}
\vspace{1.2em}

\begin{center}
\Large
$t_\mathrm{flash} = t_\mathrm{PKUP} + C$, \qquad $C \simeq -1719\ns$
\end{center}
\end{frame}

\end{document}
